Deux sphères de même rayon sont placées dans le vide. La première est un corps noir ; la seconde est un corps gris d'émissivité \num{\v{em}}. Les deux sont à \qty{\v{T}}{K}.

\begin{enumerate}[label=\alph*)]
    \item Quel est le rapport entre la puissance rayonnée par la sphère grise et celle de la sphère noire ?
    \item Les deux sphères ont-elles le même maximum d'émission $\lambda_{\max}$ ? Justifiez et calculez-le.
    \item À quelle température faudrait-il porter la sphère grise pour qu'elle rayonne la même puissance que la sphère noire à \qty{\v{T}}{K} ?
\end{enumerate}
